\documentclass[11pt]{article}

%%%%%%%%%%%%% Package imports %%%%%%%%%%%%%
\usepackage[backend=biber, style=numeric]{biblatex}
\usepackage{enumitem}
\usepackage[margin=1in]{geometry}
\usepackage{amsmath}
\usepackage{amssymb}
\usepackage{bm}
\usepackage{booktabs}
\usepackage{tabularx}
\usepackage[hidelinks]{hyperref}

\addbibresource{ekf_health_diagnostics.bib}

%%%%%%%%%%%%% Macro definitions %%%%%%%%%%%%%
\setlist{itemsep=2pt,topsep=4pt}

\newcommand{\x}{\bm{x}}
\newcommand{\xhat}{\hat{\bm{x}}}
\newcommand{\e}{\bm{e}}
\newcommand{\nuv}{\bm{\nu}}
\newcommand{\Pm}{\bm{P}}
\newcommand{\Sm}{\bm{S}}
\newcommand{\Hm}{\bm{H}}
\newcommand{\Fm}{\bm{F}}
\newcommand{\Km}{\bm{K}}
\newcommand{\Qm}{\bm{Q}}
\newcommand{\Rm}{\bm{R}}
\newcommand{\Id}{\bm{I}}
\newcommand{\Phim}{\bm{\Phi}}
\newcommand{\Jm}{\bm{J}}
\newcommand{\T}{^{\mathsf{T}}}
\newcommand{\E}{\mathbb{E}}
\newcommand{\chisq}[1]{\chi^2_{#1}}

%%%%%%%%%%%%% Title card content %%%%%%%%%%%%%
\title{Characterizing the Health and Behavior of an Extended Kalman Filter:\\
A Catalog of Diagnostic Techniques}
\author{Andrew}
\date{\today}

%%%%%%%%%%%%% Begin document %%%%%%%%%%%%%
\begin{document}
\maketitle

\begin{abstract}
This note collects techniques for assessing whether an extended Kalman filter (EKF) is
consistent, well tuned, numerically sound, and operating within the regime where its
linearizations are valid. Techniques are grouped by what they diagnose and by whether they
require knowledge of the true state (simulation only) or can be applied to real data.
\end{abstract}

\section*{Notation}
The discrete-time system is
\begin{align}
  \x_{k+1} &= \bm{f}(\x_k) + \bm{w}_k, & \bm{w}_k &\sim \mathcal{N}(\bm{0},\Qm_k),\\
  \bm{z}_k &= \bm{h}(\x_k) + \bm{v}_k, & \bm{v}_k &\sim \mathcal{N}(\bm{0},\Rm_k),
\end{align}
with Jacobians $\Fm_k = \partial \bm{f}/\partial\x|_{\xhat_k^+}$ and
$\Hm_k = \partial \bm{h}/\partial\x|_{\xhat_k^-}$. The state dimension is $n$ and the
measurement dimension is $m$. The innovation and its predicted covariance are
\begin{equation}
  \nuv_k = \bm{z}_k - \bm{h}(\xhat_k^-), \qquad
  \Sm_k = \Hm_k \Pm_k^- \Hm_k\T + \Rm_k,
\end{equation}
and the estimation error is $\e_k = \x_k - \xhat_k^+$. Monte Carlo runs are indexed by
$i = 1,\dots,N$, and time steps by $k = 1,\dots,K$. The $p$-quantile of the chi-square
distribution with $d$ degrees of freedom is written $\chisq{d}(p)$.

%======================================================================
\section{Innovation (Residual) Tests}
\label{sec:innovation}
These tests need no truth, so they are the primary diagnostics on real data.

\subsection{Normalized innovation squared (NIS)}
\begin{equation}
  \epsilon_{\nu,k} = \nuv_k\T \Sm_k^{-1} \nuv_k \;\sim\; \chisq{m}
  \quad\text{(if the filter is consistent).}
\end{equation}
For a two-sided 95\% test at a single step, accept if
$\chisq{m}(0.025) \le \epsilon_{\nu,k} \le \chisq{m}(0.975)$.
Single-step tests are noisy, so use a time average over a window of $K$ steps:
\begin{equation}
  \bar{\epsilon}_\nu = \frac{1}{K}\sum_{k=1}^{K} \epsilon_{\nu,k},
  \qquad K\bar{\epsilon}_\nu \sim \chisq{Km}.
\end{equation}
Over $N$ Monte Carlo runs, the ensemble average at each step obeys
$N\bar{\epsilon}_{\nu,k} \sim \chisq{Nm}$.

\emph{Interpretation:} A NIS that sits too high means the filter is overconfident
($\Qm$ or $\Rm$ is too small, or there is unmodeled dynamics or bias). A NIS that sits too
low means the filter is pessimistic ($\Qm$ or $\Rm$ is too large).

\subsection{Innovation whiteness}
For an optimal filter, the innovations are white. For a scalar innovation sequence (or for
each component of a vector sequence), the normalized sample autocorrelation at lag $\ell$ is
\begin{equation}
  \hat{\rho}(\ell) = \frac{\sum_{k=1}^{K-\ell} \nu_k \nu_{k+\ell}}
                          {\sum_{k=1}^{K} \nu_k^2},
  \qquad \hat{\rho}(\ell) \approx \mathcal{N}\!\left(0,\tfrac{1}{K}\right),
\end{equation}
so values outside $\pm 1.96/\sqrt{K}$ at more than about 5\% of lags indicate correlation.
For a vector sequence, the Bar-Shalom form normalizes the inner product
$\nuv_k\T\nuv_{k+\ell}$ in the same way. To test many lags jointly, use the Ljung--Box
portmanteau statistic:
\begin{equation}
  Q_{\mathrm{LB}} = K(K+2)\sum_{\ell=1}^{L}\frac{\hat{\rho}^2(\ell)}{K-\ell}
  \;\sim\; \chisq{L}.
\end{equation}
\emph{Interpretation:} Colored innovations mean the filter is extracting information
suboptimally. The usual causes are a mis-tuned $\Qm$, correlated measurement noise, or
model error.

\subsection{Zero-mean test}
The sample mean $\bar{\nuv} = K^{-1}\sum_k \nuv_k$ should satisfy
$K\,\bar{\nuv}\T \bar{\Sm}^{-1} \bar{\nuv} \sim \chisq{m}$, approximately, where
$\bar{\Sm}$ is the average $\Sm_k$. A nonzero mean points to an unmodeled sensor bias,
a frame misalignment, or a timing offset.

\subsection{Innovation envelopes and residual ratios}
\begin{itemize}
  \item Plot each innovation component against $\pm 3\sqrt{[\Sm_k]_{jj}}$.
  \item Compare pre-fit residuals $\bm{z}_k - \bm{h}(\xhat_k^-)$ with post-fit residuals
        $\bm{z}_k - \bm{h}(\xhat_k^+)$. Their ratio shows how strongly each update pulls
        the state.
\end{itemize}

%======================================================================
\section{Truth-Based Consistency (Simulation and Monte Carlo)}
\label{sec:truth}

\subsection{Normalized estimation error squared (NEES)}
\begin{equation}
  \epsilon_k = \e_k\T (\Pm_k^+)^{-1} \e_k \;\sim\; \chisq{n}.
\end{equation}
The average NEES (ANEES) over $N$ runs is
\begin{equation}
  \bar{\epsilon}_k = \frac{1}{N}\sum_{i=1}^{N}\epsilon_k^{(i)},
  \qquad N\bar{\epsilon}_k \sim \chisq{Nn},
\end{equation}
with the 95\% acceptance region
$\left[\chisq{Nn}(0.025)/N,\; \chisq{Nn}(0.975)/N\right]$.
NEES is the most direct check of whether $\Pm$ honestly describes the actual error.

\subsection{Per-state error envelopes}
Plot $e_{k,j}$ against $\pm 3\sqrt{[\Pm_k^+]_{jj}}$ for every state $j$. The scalar NEES
can hide a single inconsistent state, and these plots show which state is the problem.
It is also useful to compute NEES separately for meaningful subsystem blocks of the state.

\subsection{RMSE versus predicted standard deviation}
\begin{equation}
  \mathrm{RMSE}_{k,j} = \sqrt{\frac{1}{N}\sum_{i=1}^{N}\big(e_{k,j}^{(i)}\big)^2}
  \quad\text{vs.}\quad
  \overline{\sigma}_{k,j} = \sqrt{\frac{1}{N}\sum_{i=1}^{N}[\Pm_k^{+(i)}]_{jj}}.
\end{equation}
An RMSE persistently above $\overline{\sigma}$ means the filter is overconfident
(``smug'').

\subsection{Ensemble bias}
The ensemble mean error $\bar{\e}_k = N^{-1}\sum_i \e_k^{(i)}$ should be statistically
indistinguishable from zero. Test it against $\Pm_k^+/N$.

\subsection{Filter--smoother comparison}
Run a Rauch--Tung--Striebel (RTS) smoother and examine $\xhat_{k|K} - \xhat_{k|k}$.
Large or structured differences reveal intervals where the forward filter struggles,
for example during observability gaps.

%======================================================================
\section{Covariance Health and Numerics}
\label{sec:numerics}
\begin{itemize}
  \item \textbf{Symmetry:} monitor $\|\Pm - \Pm\T\|_F / \|\Pm\|_F$.
  \item \textbf{Positive definiteness:} monitor $\lambda_{\min}(\Pm)$ and any
        Cholesky failures.
  \item \textbf{Conditioning:} monitor $\kappa(\Pm) = \lambda_{\max}/\lambda_{\min}$.
        Growth warns of numerical trouble or of nearly unobservable directions.
  \item \textbf{Collapse and divergence:} track $\operatorname{tr}\Pm_k$ and
        $\log\det\Pm_k$ over time. A collapsing $\Pm$ drives the gains toward zero, so the
        filter stops listening to measurements.
  \item \textbf{Correlation coefficients:}
        $r_{ij} = P_{ij}/\sqrt{P_{ii}P_{jj}} \in [-1,1]$. Values pinned near $\pm1$
        indicate states that cannot be separated.
  \item \textbf{Update-form robustness:} compare the standard update with the Joseph form,
        \begin{equation}
          \Pm^+ = (\Id - \Km\Hm)\Pm^-(\Id - \Km\Hm)\T + \Km\Rm\Km\T,
        \end{equation}
        and with square-root or $UDU\T$ factorized implementations. If the results
        differ, the problem is numerical rather than a modeling issue.
\end{itemize}

%======================================================================
\section{Linearization Validity}
\label{sec:linearization}
Linearization error is the failure mode specific to the EKF.

\subsection{Jacobian verification}
Check the analytic $\Fm$ and $\Hm$ against complex-step derivatives, which have no
subtractive cancellation error:
\begin{equation}
  \frac{\partial f}{\partial x_j} \approx
  \frac{\operatorname{Im}\, f(\x + i h\,\bm{\hat{e}}_j)}{h},
  \qquad h \sim 10^{-20},
\end{equation}
or against central finite differences when $f$ is not complex-analytic.

\subsection{Nonlinearity index}
Following Junkins and Singla, sample points $\x_s$ on the $1\sigma$--$3\sigma$ covariance
ellipsoid about $\bar{\x}$ and compute
\begin{equation}
  \nu(t) = \max_s
  \frac{\|\Phim(t;\x_s) - \Phim(t;\bar{\x})\|}{\|\Phim(t;\bar{\x})\|}.
\end{equation}
A large $\nu$ over a propagation interval means the linear covariance propagation cannot be
trusted over that interval. An analogous index applies to $\Hm$.

\subsection{Covariance propagation cross-check}
Propagate the same Gaussian prior through Monte Carlo sampling or the unscented
transform, and compare the result with the EKF prediction
$(\bar{\x}_{\mathrm{EKF}},\Pm_{\mathrm{EKF}})$. A scalar summary is the KL divergence
between the two Gaussians:
\begin{equation}
  D_{\mathrm{KL}} = \tfrac12\!\left[
  \operatorname{tr}\!\big(\Pm_{\mathrm{EKF}}^{-1}\Pm_{\mathrm{ref}}\big)
  + \Delta\bar{\x}\T\Pm_{\mathrm{EKF}}^{-1}\Delta\bar{\x}
  - n + \ln\frac{\det\Pm_{\mathrm{EKF}}}{\det\Pm_{\mathrm{ref}}}\right].
\end{equation}

\subsection{Second-order term magnitude}
Estimate the second-order mean correction
$\tfrac12\sum_j \bm{\hat{e}}_j \operatorname{tr}(\nabla^2 f_j\,\Pm)$ and compare it with
$\sqrt{\operatorname{diag}\Pm}$. If the correction is comparable in size, a second-order
filter may be warranted.

\subsection{Alternative-filter comparison}
Run an iterated EKF or an unscented Kalman filter on the same data. Meaningful divergence
between their estimates indicates that linearization error is material for the problem.

\subsection{Iterated measurement update}
\label{sec:iekf}
When $\Hm$ varies appreciably over the prior uncertainty, relinearize the update about the
latest iterate $\bm{\eta}_i$, starting at $\bm{\eta}_1=\bar{\x}$ (Gauss--Newton on the
MAP cost):
\begin{align}
  \bm{\nu}_i &= \bm{z} - h(\bm{\eta}_i) - \Hm_i(\bar{\x} - \bm{\eta}_i), \\
  \Sm_i &= \Hm_i\Pm^-\Hm_i\T + \Rm, \qquad
  \Km_i = \Pm^-\Hm_i\T\Sm_i^{-1}, \\
  \bm{\eta}_{i+1} &= \bar{\x} + \Km_i\bm{\nu}_i .
\end{align}
Iteration stops when
$(\bm{\eta}_{i+1}-\bm{\eta}_i)\T(\Pm^-)^{-1}(\bm{\eta}_{i+1}-\bm{\eta}_i)<10^{-6}$ or at
a maximum count. The posterior mean is the last iterate, and the Joseph-form covariance
uses the final $\Hm$ and $\Km$.

The health statistics keep their definitions at the prior. The NIS uses the first-pass
$\bm{\nu}_1$ and $\Sm_1$, which are the innovation and covariance about $\bar{\x}$, so it
remains $\chi^2_m$ for a consistent filter, while the NEES uses the final posterior. In
the tethered three-body demo, a single linearization of the payload--catch-mechanism
crosslink range over a 20\,km position uncertainty at 7500\,km separation leaves a
linearization error of about $200\sigma$ in that row; this alone drove the first-update
ANEES to $\sim\!10^3$ times its expected value, and five iterations restore it to $\sim\!1$.

%======================================================================
\section{Observability and Information}
\label{sec:observability}

\subsection{Local observability}
Use the observability matrix
$\mathcal{O}_k = [\Hm_k\T,\ (\Hm_{k+1}\Fm_k)\T,\ \dots]\T$, or the information Gramian
\begin{equation}
  \bm{W}_k = \sum_{j=k}^{k+L}\Phim_{j,k}\T \Hm_j\T \Rm_j^{-1}\Hm_j \Phim_{j,k}.
\end{equation}
Track its rank, its smallest singular value, and its condition number along the trajectory.
The smallest singular value serves as a degree of observability.

\subsection{Per-measurement information}
Measure each update's contribution through the eigenvalues of $\Pm^-$ versus $\Pm^+$, or
through $\Delta(\Pm^{-1}) = \Hm\T\Rm^{-1}\Hm$, or through the entropy reduction
$\tfrac12\ln(\det\Pm^-/\det\Pm^+)$. Measurements that contribute almost nothing deserve
investigation.

\subsection{Posterior Cram\'er--Rao lower bound (PCRLB)}
For additive Gaussian noise, the Tichavsk\'y et al.\ recursion for the information
matrix is
\begin{equation}
  \Jm_{k+1} = \Qm_k^{-1} + \E\!\left[\Hm_{k+1}\T\Rm_{k+1}^{-1}\Hm_{k+1}\right]
  - \Qm_k^{-1}\E[\Fm_k]\left(\Jm_k + \E\!\left[\Fm_k\T\Qm_k^{-1}\Fm_k\right]\right)^{-1}
  \E[\Fm_k]\T\Qm_k^{-1},
\end{equation}
with expectations evaluated along the true trajectories. Compare $\Jm_k^{-1}$ with both
$\Pm_k$ and the Monte Carlo mean-squared error. This separates filter inefficiency from
the inherent difficulty of the problem.

%======================================================================
\section{Tuning and Adaptive Diagnostics}
\label{sec:tuning}
\begin{itemize}
  \item \textbf{Covariance matching:} over a window of $W$ steps, compute
        $\hat{\bm{C}}_\nu = W^{-1}\sum \nuv_k\nuv_k\T$ and compare it with $\Sm_k$.
        This suggests the estimate $\hat{\Rm} = \hat{\bm{C}}_\nu - \Hm\Pm^-\Hm\T$.
  \item \textbf{Mehra's method and autocovariance least squares (ALS):} these estimate
        $\Qm$ and $\Rm$ from the innovation autocorrelations of a suboptimal filter,
        which gives a data-driven tuning check.
  \item \textbf{Gain history:} plot $\|\Km_k\|$ over time. Gains decaying toward zero
        while errors persist are the classic sign that $\Qm$ is too small.
  \item \textbf{Consider and sensitivity analysis:} propagate the sensitivity
        $\partial\xhat/\partial\bm{c}$ to unestimated parameters $\bm{c}$ (for example
        sensor biases or dynamics coefficients) and form the consider covariance.
        This tests whether such parameters explain an observed inconsistency.
  \item \textbf{Error budget:} enable each error source individually in Monte Carlo runs
        to obtain per-source contributions to the total error.
\end{itemize}

%======================================================================
\section{Fault and Outlier Detection}
\label{sec:fault}
\begin{itemize}
  \item \textbf{Chi-square gating:} reject $\bm{z}_k$ if
        $\epsilon_{\nu,k} > \chisq{m}(1-\alpha)$. Track the rejection rate over time;
        a rising rate signals model mismatch rather than random outliers.
  \item \textbf{CUSUM on NIS:} the statistic
        \begin{equation}
          g_k = \max\!\big(0,\; g_{k-1} + \epsilon_{\nu,k} - m - \delta\big)
        \end{equation}
        raises an alarm when $g_k > h$. It detects slow drifts that per-step tests miss.
        The sequential probability ratio test (SPRT) is a related alternative.
  \item \textbf{Generalized likelihood ratio (GLR):} hypothesizes a fault signature and
        estimates its onset time and magnitude. It is suited to step faults and jumps.
\end{itemize}

%======================================================================
\section{Stability Theory}
\label{sec:stability}
Reif et al.\ give sufficient conditions for the EKF estimation error to be exponentially
bounded in mean square. The conditions require uniform bounds on $\Fm_k$, $\Hm_k$, and
$\Pm_k$, uniform observability, and sufficiently small initial error and noise. They are
more useful for proving boundedness than for online monitoring.

%======================================================================
\section{Attitude-Specific Considerations}
\label{sec:attitude}
\begin{itemize}
  \item \textbf{Additive quaternion EKF:} monitor $\big|\|\bm{q}\| - 1\big|$, and note
        that the $4\times4$ covariance is singular along the norm direction.
  \item \textbf{Multiplicative EKF (MEKF):} verify the covariance reset after each update,
        and confirm that the attitude error stays small enough for the small-angle
        error parameterization to hold.
  \item \textbf{Coupled orbit--attitude filters:} compute NEES separately for the orbit
        block, the attitude block, and the cross-covariance. A whole-state NEES can look
        acceptable while one subsystem is badly inconsistent.
\end{itemize}

%======================================================================
\section{Summary}
\begin{center}
\small
\begin{tabularx}{\textwidth}{@{}lcX@{}}
\toprule
\textbf{Technique} & \textbf{Truth?} & \textbf{Primarily detects} \\
\midrule
NIS (single step, time-averaged) & No  & Over- or under-confidence; mis-tuned $\Qm$, $\Rm$ \\
Innovation whiteness, Ljung--Box & No  & Suboptimal gain, correlated noise, model error \\
Innovation zero-mean             & No  & Biases, misalignment, timing offsets \\
NEES and ANEES                   & Yes & Honesty of $\Pm$ \\
Per-state error envelopes        & Yes & Which states are inconsistent \\
RMSE vs.\ $\overline{\sigma}$    & Yes & Filter smugness \\
Filter--smoother difference      & No  & Struggling intervals, observability gaps \\
Covariance symmetry, PD, $\kappa$& No  & Numerical breakdown \\
$\operatorname{tr}\Pm$, $\log\det\Pm$ & No & Collapse or divergence \\
Jacobian check (complex step)    & No  & Implementation bugs \\
Nonlinearity index, UT/MC check  & No  & Linearization breakdown \\
IEKF or UKF comparison           & No  & Material linearization error \\
Observability Gramian            & No  & Weakly observable directions \\
PCRLB                            & Yes & Filter inefficiency vs.\ problem limits \\
Covariance matching, ALS         & No  & Noise covariance estimation \\
Gain history                     & No  & $\Qm$ too small \\
Consider analysis, error budget  & Mixed & Unmodeled parameters, error sources \\
Gating rate, CUSUM, GLR          & No  & Outliers, drifts, faults \\
\bottomrule
\end{tabularx}
\end{center}

%%%%%%%%%%%%% Bibliography %%%%%%%%%%%%%
\nocite{*}
\printbibliography

\end{document}
